\chapter{Pruebas y Validación del Módulo de Clasificación}
\label{chap:pruebas}

Este capítulo describe la \textbf{metodología y el diseño experimental} empleados para la evaluación del rendimiento del Módulo de Clasificación de Secuencias de LSM. La validación se centra en establecer el protocolo para probar la precisión del modelo de \textit{Deep Learning} (LSTM) en las 29 palabras objetivo, utilizando el aplicativo web \textbf{Flask} y el servicio \texttt{trainer}.

\section{Metodología de Pruebas: Adquisición y Diseño del \textit{Test Set}}

La prueba del sistema se fundamenta en la adquisición de un \textbf{Conjunto de Pruebas (\textit{Test Set})} independiente, capturado mediante el aplicativo web Flask para garantizar la fidelidad de los datos.

\subsection{Protocolo de Captura de Secuencias de Prueba}

La adquisición de las muestras de prueba se ejecuta a través de la interfaz web, donde el \textit{frontend} utiliza \textbf{MediaPipe Tasks Vision} para la extracción de características en tiempo real.

\begin{enumerate}
    \item \textbf{Adquisición de \textit{Landmarks}:} El usuario ejecuta una de las 29 señas frente a la cámara. La aplicación registra una secuencia de \textit{frames}, donde cada \textit{frame} contiene:
    \begin{itemize}
        \item \textbf{Pose (17 puntos filtrados):} Coordenadas suavizadas ($X, Y$), y coordenadas originales de profundidad y visibilidad ($Z, V$). Las coordenadas $X, Y$ son sometidas a un filtro \textbf{Exponencial Móvil Promedio (EMA)} para mitigar el \textit{jitter} temporal.
        \item \textbf{Manos (42 puntos):} Coordenadas ($X, Y, Z$) para los 21 \textit{landmarks} de cada mano.
    \end{itemize}
    \item \textbf{Criterio de Cierre de Muestra:} La captura finaliza automáticamente al detectar un periodo prolongado de estabilidad, asegurando que solo se registren secuencias completas y bien ejecutadas de la seña, lo que garantiza la calidad del conjunto de validación.
    \item \textbf{Almacenamiento Persistente:} La secuencia de \textit{landmarks} se transfiere al \textit{backend} (Flask). Este la almacena como un \textit{dataset} en un archivo \textbf{HDF5} (\texttt{.h5}) y registra la \textit{metadata} (\textit{Ground Truth}, duración, y número de \textit{frames}) en \textbf{PostgreSQL}.
\end{enumerate}

\subsection{Diseño del Conjunto de Validación por Clase}

Para una validación rigurosa, se diseñó el \textit{Test Set} bajo los siguientes parámetros de control experimental. La Tabla \ref{tab:diseño_test_set} lista las 29 palabras que componen el vocabulario de prueba y los requisitos de captura.

\begin{itemize}
    \item \textbf{Independencia:} Las secuencias de prueba deben ser grabadas por un grupo de colaboradores \textbf{distinto} al utilizado para generar el conjunto de datos de entrenamiento.
    \item \textbf{Cobertura Mínima:} Se establece una cantidad mínima de \textbf{$M$ secuencias} por cada una de las 29 palabras objetivo.
\end{itemize}

\begin{table}[h!]
    \centering
    \caption{Diseño del Conjunto de Validación: Vocabulario y Criterios de Muestra.}
    \label{tab:diseño_test_set}
    \begin{tabular}{|c|l|c|c|}
        \hline
        \textbf{ID} & \textbf{Palabra (LSM)} & \textbf{Muestras Mínimas por Seña ($M$)} & \textbf{Tipo de Muestra} \\
        \hline
        01 & inyección & $M$ & Secuencia Temporal (HDF5) \\
        02 & agua & $M$ & Secuencia Temporal (HDF5) \\
        03 & farmacia & $M$ & Secuencia Temporal (HDF5) \\
        04 & sida & $M$ & Secuencia Temporal (HDF5) \\
        05 & salud & $M$ & Secuencia Temporal (HDF5) \\
        06 & curación & $M$ & Secuencia Temporal (HDF5) \\
        07 & virus & $M$ & Secuencia Temporal (HDF5) \\
        08 & operación & $M$ & Secuencia Temporal (HDF5) \\
        09 & enfermo & $M$ & Secuencia Temporal (HDF5) \\
        10 & temperatura & $M$ & Secuencia Temporal (HDF5) \\
        11 & cómo & $M$ & Secuencia Temporal (HDF5) \\
        12 & orinar & $M$ & Secuencia Temporal (HDF5) \\
        13 & por favor & $M$ & Secuencia Temporal (HDF5) \\
        14 & limpiar & $M$ & Secuencia Temporal (HDF5) \\
        15 & hola & $M$ & Secuencia Temporal (HDF5) \\
        16 & síntoma & $M$ & Secuencia Temporal (HDF5) \\
        17 & infección & $M$ & Secuencia Temporal (HDF5) \\
        18 & cita & $M$ & Secuencia Temporal (HDF5) \\
        19 & doctor & $M$ & Secuencia Temporal (HDF5) \\
        20 & bien & $M$ & Secuencia Temporal (HDF5) \\
        21 & cáncer & $M$ & Secuencia Temporal (HDF5) \\
        22 & menstruación & $M$ & Secuencia Temporal (HDF5) \\
        23 & influenza & $M$ & Secuencia Temporal (HDF5) \\
        24 & mal & $M$ & Secuencia Temporal (HDF5) \\
        25 & levantar & $M$ & Secuencia Temporal (HDF5) \\
        26 & emergencia & $M$ & Secuencia Temporal (HDF5) \\
        27 & recostar & $M$ & Secuencia Temporal (HDF5) \\
        28 & radio & $M$ & Secuencia Temporal (HDF5) \\
        29 & diarrea & $M$ & Secuencia Temporal (HDF5) \\
        \hline
    \end{tabular}
\end{table}


\section{Traducción y Métrica de Clasificación de Secuencias}

Este proceso describe la transformación de los datos crudos (\textit{landmarks}) en la entrada del modelo y la definición de las métricas que serán utilizadas en el Capítulo de Resultados.

\subsection{Preprocesamiento y Entrada al Modelo LSTM}

El servicio \texttt{trainer} procesa las secuencias para adaptarlas al formato de entrada del modelo LSTM, que es crucial para la clasificación de series de tiempo:

\begin{enumerate}
    \item \textbf{Normalización Temporal:} Se añade una característica clave a cada \textit{frame}: el \textbf{progreso relativo de tiempo} de la seña. Esta característica, normalizada entre $[0, 1]$, permite que el modelo ignore las variaciones en la velocidad de ejecución, esencial para la invariancia del sistema.
    \begin{equation}
        \text{Progreso}_t = \frac{\text{Timestamp}_t - \text{Timestamp}_{\text{inicio}}}{\text{Timestamp}_{\text{fin}} - \text{Timestamp}_{\text{inicio}}}
    \end{equation}
    \item \textbf{Alineación (\textit{Padding}):} Se aplica \textit{Zero-Padding} a todas las secuencias para alcanzar la longitud máxima $T_{\text{max}}$, estandarizando la forma de entrada a la red a $$(N_{\text{muestras}}, T_{\text{max}}, F_{\text{total}})$$.
\end{enumerate}

\subsection{Inferencia, Traducción (*Argmax*) y Métricas}

\begin{enumerate}
    \item \textbf{Inferencia y Traducción:} El modelo LSTM produce un vector de probabilidad $\mathbf{P} \in \mathbb{R}^{29}$. La decisión de clasificación se obtiene aplicando la función \textit{Argmax}, seleccionando la clase con la probabilidad más alta:
    \begin{equation}
        \text{Clase Predicha} = \underset{i}{\operatorname{argmax}} (P_i)
    \end{equation}
    Este resultado se interpreta a través de un \textbf{mapeo uno a uno} (1:1) de la clase numérica a la etiqueta de texto correspondiente. Por lo tanto, el sistema \textbf{reconoce} la seña del lenguaje de señas mexicano (LSM) como una de las 29 palabras objetivo.

    \textbf{Mapeo a Audio (JavaScript):} Una vez que el sistema \textbf{reconoce} la palabra, el \textit{frontend} utiliza la etiqueta de texto resultante para generar la salida de audio mediante una función de Texto a Voz (TTS), completando así la traducción.
    \begin{verbatim}
// Mapa de la clase numérica a la etiqueta de palabra (simulación)
const LABEL_MAP = { 0: 'inyección', 1: 'agua', /* ... */ 28: 'diarrea' };

function playAudio(predictedClassIndex) {
    // 1. Mapeo 1:1 de la clase numérica a la palabra de texto
    const recognizedWord = LABEL_MAP[predictedClassIndex];

    // 2. Uso de la API de Speech Synthesis (TTS) del navegador
    const utterance = new SpeechSynthesisUtterance(recognizedWord);

    // 3. Reproducción del sonido de la palabra reconocida
    window.speechSynthesis.speak(utterance);

    return recognizedWord;
}
    \end{verbatim}

    \item \textbf{Métrica de Evaluación:} Se utilizará el \textbf{F1-Score} como métrica principal, ya que proporciona un indicador equilibrado del rendimiento del clasificador en todas las clases, siendo robusto frente a desequilibrios entre la Precisión ($P$) y el Recall ($R$):
    \begin{equation}
        F1 = 2 \cdot \frac{P \cdot R}{P + R}
    \end{equation}
\end{enumerate}

El cálculo de estas métricas sobre el \textit{Test Set} será presentado y analizado en el Capítulo de Resultados.